\documentclass[10pt,a4paper,landscape]{article}
\usepackage[utf8]{inputenc}
\usepackage[T1]{fontenc}
\usepackage[french]{babel}
\frenchbsetup{StandardLists=true}
\usepackage[landscape,left=1cm,right=1cm,top=1.5cm,bottom=1.5cm]{geometry}
\usepackage{multirow,tabularx,booktabs}
\usepackage[dvipsnames]{xcolor}
\usepackage{fouriernc}
\usepackage{fancyhdr}
\usepackage{colortbl}
\usepackage{multicol}
\setlength{\columnsep}{8mm}
\setlength{\columnseprule}{0.4pt}
\pagestyle{fancy}
\renewcommand\headrulewidth{1pt}
\fancyhead[L]{Lycée Denis Diderot}
\fancyhead[R]{Classe de 3\textsuperscript{ème}}
\fancyfoot[C]{}
\fancyfoot[R]{Année 2025-2026}
\fancyfoot[L]{Alexis RUIZ}
\setlength{\parindent}{0pt}
\renewcommand{\arraystretch}{1.35}

% Couleurs
\definecolor{exoheader}{RGB}{30,90,160}
\definecolor{exolight}{RGB}{220,235,255}
\definecolor{totalrow}{RGB}{255,230,200}

% Macro en-tête de partie
\newcommand{\exotitle}[2]{%
  \noindent\colorbox{exoheader}{%
    \parbox{\dimexpr\linewidth-2\fboxsep\relax}{%
      \textcolor{white}{\textbf{\large #1 \hfill #2}}%
    }%
  }%
}

\begin{document}

\begin{center}
  {\LARGE\textbf{GRILLE DE CORRECTION}}\\[1mm]
  {\large Activité documentaire -- Introduction à l'électricité \hspace{1cm} 3\textsuperscript{ème} \hspace{1cm} \textbf{Total : /20 points}}
\end{center}

\vspace{2mm}

\begin{multicols}{2}

% ===================== DOCUMENT 1 =====================
\exotitle{Document 1 --- L'électricité dans notre quotidien}{/6 pts}

\vspace{1mm}

\begin{tabularx}{\linewidth}{|c|X|c|}
  \hline
  \rowcolor{exolight}
  \textbf{Q} & \textbf{Réponse attendue} & \textbf{Pts} \\
  \hline
  \multirow{3}{*}{1} & 1\textsuperscript{er} domaine (au choix : éclairage, électroménager, transports, communication, médecine) & 1 \\
  \cline{2-3}
  & 2\textsuperscript{e} domaine & 1 \\
  \cline{2-3}
  & 3\textsuperscript{e} domaine & 1 \\
  \hline
  2 & Centrales \textbf{nucléaires} -- environ \textbf{70\,\%} & 1 \\
  \hline
  \multirow{2}{*}{3} & Tension dans les foyers : \textbf{230\,V} & 1 \\
  \cline{2-3}
  & Transport par \textbf{lignes à haute tension} (400\,000\,V), puis transformée & 1 \\
  \hline
  \rowcolor{totalrow}
  \multicolumn{2}{|r|}{\textbf{Sous-total Document 1}} & \textbf{/6} \\
  \hline
\end{tabularx}

\vspace{4mm}

% ===================== DOCUMENT 2 =====================
\exotitle{Document 2 --- Les composants d'un circuit}{/7 pts}

\vspace{1mm}

\begin{tabularx}{\linewidth}{|c|X|c|}
  \hline
  \rowcolor{exolight}
  \textbf{Q} & \textbf{Réponse attendue} & \textbf{Pts} \\
  \hline
  \multirow{2}{*}{4} & Rôle : \textbf{fournit l'énergie électrique} au circuit & 1 \\
  \cline{2-3}
  & 2 exemples parmi : pile, batterie, panneau solaire, dynamo (0,5 pt chacun) & 1 \\
  \hline
  5 & \textit{(tableau)} Récepteur : convertit l'énergie, ex. lampe/moteur & 1 \\
  \cline{2-3}
  & Interrupteur : ouvre/ferme le circuit, ex. bouton ON/OFF & 1 \\
  \cline{2-3}
  & Fils conducteurs : circulation du courant, ex. fils de cuivre & 1 \\
  \hline
  \multirow{2}{*}{6} & Circuit \textbf{fermé} : courant circule $\Rightarrow$ lampe \textbf{allumée} & 1 \\
  \cline{2-3}
  & Circuit \textbf{ouvert} : courant ne circule plus $\Rightarrow$ lampe \textbf{éteinte} & 1 \\
  \hline
  \rowcolor{totalrow}
  \multicolumn{2}{|r|}{\textbf{Sous-total Document 2}} & \textbf{/7} \\
  \hline
\end{tabularx}

\vspace{4mm}

% ===================== DOCUMENT 3 =====================
\exotitle{Document 3 --- Dangers et sécurité électrique}{/5 pts}

\vspace{1mm}

{\footnotesize\textit{Q7 : 1 pt par danger correctement nommé et expliqué. Q8 : 1 pt par règle.}}

\vspace{1mm}

\begin{tabularx}{\linewidth}{|c|X|c|}
  \hline
  \rowcolor{exolight}
  \textbf{Q} & \textbf{Réponse attendue} & \textbf{Pts} \\
  \hline
  \multirow{2}{*}{7} & 1\textsuperscript{er} danger + explication (électrocution, incendie ou brûlures) & 1 \\
  \cline{2-3}
  & 2\textsuperscript{e} danger + explication & 1 \\
  \hline
  \multirow{3}{*}{8} & 1\textsuperscript{re} règle (au choix parmi les 5 du document) & 1 \\
  \cline{2-3}
  & 2\textsuperscript{e} règle & 1 \\
  \cline{2-3}
  & 3\textsuperscript{e} règle & 1 \\
  \hline
  \rowcolor{totalrow}
  \multicolumn{2}{|r|}{\textbf{Sous-total Document 3}} & \textbf{/5} \\
  \hline
\end{tabularx}

\vspace{4mm}

% ===================== SYNTHÈSE =====================
\exotitle{Questions de synthèse}{/2 pts}

\vspace{1mm}

\begin{tabularx}{\linewidth}{|c|X|c|}
  \hline
  \rowcolor{exolight}
  \textbf{Q} & \textbf{Réponse attendue} & \textbf{Pts} \\
  \hline
  9 & L'eau conduit le courant. À 230\,V (> 25\,V) $\Rightarrow$ risque d'\textbf{électrocution} mortelle & 1 \\
  \hline
  10 & 4,5\,V $<$ 25\,V (seuil de danger) $\Rightarrow$ \textbf{aucun danger} pour l'élève & 1 \\
  \hline
  \rowcolor{totalrow}
  \multicolumn{2}{|r|}{\textbf{Sous-total Synthèse}} & \textbf{/2} \\
  \hline
\end{tabularx}

\vspace{6mm}

\noindent\colorbox{exoheader}{%
  \parbox{\dimexpr\linewidth-2\fboxsep\relax}{%
    \textcolor{white}{\textbf{\Large TOTAL \hfill /20 points}}%
  }%
}

\end{multicols}

\end{document}
